% !PRINTED 509
\weberpage{509}

To prove it, we denote by $t$ a variable which does not occur in the $\alpha$ and in $\omega$, and set

\weq{11}{\varphi(t) = t^m + \alpha_1 t^{m-1} + \cdots + \alpha_m.}

Then, if $n$ is the degree of the field $\Omega$,
\[
\Phi(t) = N[\varphi(t)]
\]
is an integral function of $t$ of degree $mn$, whose coefficients are integral rational functionals. At the same time $\Phi(\omega) = 0$, and consequently $\omega$ is an integer.

\gesperrt{From this it follows further that Theorem 6. remains correct if the coefficients $A_{k,i}$ are integral functionals not in $R$, but in $\Omega$.}

If an integral functional $\omega$ is at the same time a number, then it is an integer, because in this case $N(t-\omega)$ has integral rational numbers as coefficients. The integers which we considered in \S.\,133 are accordingly contained as special cases among the integral functionals.

\S.\,138.

\gesperrt{Divisibility. Associated functionals. Units.}

The integral functionals are subject to laws of divisibility similar to those of the integral rational numbers. To recognise them, we place the following definition at the head:

1. \gesperrt{An integral functional $\alpha$ is called divisible by another integral functional $\beta$, different from zero, if the quotient $\alpha:\beta=\gamma$ is an integral functional.}

Then $\alpha=\beta\gamma$, and $\beta$ and $\gamma$ are called \gesperrt{divisors} of $\alpha$, and one also says that \gesperrt{$\beta$ and $\gamma$ divide $\alpha$}. The number 0 is divisible by \gesperrt{every integral} functional.

From this definition the following fundamental theorems on divisibility are obtained without difficulty:

2. \gesperrt{If $\alpha$ and $\alpha_1$ are divisible by $\beta$, then $\alpha\pm\alpha_1$ is also divisible by $\beta$.}

For if $\alpha=\beta\gamma$, $\alpha_1=\beta\gamma_1$, then $\alpha\pm\alpha_1=\beta(\gamma\pm\gamma_1)$, and if $\gamma$ and $\gamma_1$ are integral, then $\gamma\pm\gamma_1$ is also integral. % [sic] source prints ``nnd'' for ``und''
